% AST2000 Part 7, challenge D: exploring the star (exam-style answers, a separate hand-in).
% Compile: cd "Project 7/paper" && latexmk -pdf part7_star.tex   (or: python build_papers.py 7 from the repo root)
% Every number written as \val{key} comes from results.tex, made by star.py (pro_7.py).
\documentclass[11pt,a4paper]{article}
\usepackage[margin=2.3cm]{geometry}
\usepackage{../../common/paper/ast2000}
\setlength{\parskip}{0.5em}
\setlength{\parindent}{0pt}

\title{AST2000 Part 7D: Exploring Our Star}
\author{Sakarias Frette}
\date{\today}

\begin{document}
\maketitle

Solar system seed \val{seed}.
Our star has mass $M=\val{star-M}\,\Msun$, radius $R=\val{star-R}\,R_\odot$ ($\val{star-R-km}$\,km) and surface temperature $T=\val{star-T}$\,K.
The numbers below are computed by \code{Project 7/star.py}.

% =====================================================================
\section*{1. Is our star a main-sequence star?}
A black body of radius $R$ and surface temperature $T$ radiates $\sigma T^4$ per unit area from its whole surface:
\begin{equation}
  L=4\pi R^2\sigma T^4=\val{star-L}\,\mathrm{W}=\val{star-L-sun}\,L_\odot .
\end{equation}
In an HR diagram (\cref{fig:hr}) our star sits on the main sequence, right next to the Sun.
That fits its mass: $\val{star-M}\,\Msun$ is slightly more than the Sun, so on the main sequence it should be slightly hotter and brighter, as it is.
It is a main-sequence star, burning hydrogen in its core.

\begin{figure}[h]\centering\safegraphic[width=0.62\textwidth]{d1_hr_diagram}
\caption{HR diagram: 50\,000 stars from \code{ast2000tools}' \code{StarPopulation}, our star and the Sun.}\label{fig:hr}\end{figure}

% =====================================================================
\section*{2. Core temperature}
\textbf{Assumptions:}
\begin{itemize}
\item uniform density $\rho_0=M/(\tfrac43\pi R^3)=\val{star-rho0}$\,kg/m$^3$;
\item hydrostatic equilibrium;
\item an ideal gas;
\item composition 75\% hydrogen and 25\% helium by mass.
\end{itemize}

\subsection*{(a) The mass profile}
With uniform density, the mass inside radius $r$ is the density times the volume of that sphere:
\begin{equation}
  M(r)=\int_0^r\rho_0\,4\pi r'^2\,\dd r'=\frac{4\pi}{3}\rho_0r^3=M\frac{r^3}{R^3}.
\end{equation}

\subsection*{(b) The temperature gradient}
\textit{Hydrostatic equilibrium.}
A shell of gas does not move: the pressure difference across it balances the gravity of the mass inside it (only the mass inside pulls, by the shell theorem):
\begin{equation}
  \frac{\dd P}{\dd r}=-\frac{GM(r)\rho_0}{r^2}=-\frac{4\pi}{3}G\rho_0^2\,r .
\end{equation}

\textit{Ideal gas.}
$P=nkT=\rho_0kT/(\mu m_H)$.
Since $\rho_0$ is constant, $\dd P/\dd r=(\rho_0k/\mu m_H)\,\dd T/\dd r$.
Equating the two expressions gives
\begin{equation}
  \frac{\dd T}{\dd r}=-\frac{4\pi}{3}G\rho_0\frac{\mu m_H}{k}\,r .
\end{equation}

\textit{Mean molecular weight.}
In the hot core all atoms are fully ionised.
Per unit mass there are $X/m_H$ hydrogen nuclei, each giving 2 particles (proton and electron), and $Y/4m_H$ helium nuclei, each giving 3 particles (nucleus and 2 electrons).
So $1/\mu=2X+\tfrac34Y$, which gives $\mu=\val{star-mu}$.

\subsection*{(c) The core temperature}
Integrating from the centre to the surface:
\begin{equation}
  T(R)-T_c=-\frac{4\pi}{3}G\rho_0\frac{\mu m_H}{k}\int_0^Rr\,\dd r=-\frac{2\pi}{3}G\rho_0\frac{\mu m_H}{k}R^2,
\end{equation}
so
\begin{equation}
  T_c=T(R)+\frac{2\pi}{3}G\rho_0\frac{\mu m_H}{k}R^2=\val{star-Tc}\times10^6\,\mathrm{K}.
\end{equation}
The temperature falls as a parabola from the centre to the surface (\cref{fig:core}, left).

If the gas were \emph{neutral} ($1/\mu=X+Y/4$, $\mu=\val{star-mu-neutral}$), $T_c$ would double to \val{star-Tc-neutral}\,MK, close to the real Sun's 15.7\,MK.
That agreement is accidental: the uniform-density model underestimates the core temperature, and the wrong $\mu$ happens to compensate.

% =====================================================================
\section*{3. Energy production and luminosity}
Since $T_c=\val{star-Tc}$\,MK is far below 90\,MK, the core fuses hydrogen through the pp chain and the CNO cycle (\val{star-branch}).
The core composition is $X=\val{star-X}$, $Y=\val{star-Y}$, $X_{\rm CNO}=\val{star-Xcno}$.
Following the AST2000 lecture notes, the energy released per kg of core is approximated by power laws in temperature ($T_6=T/10^6$\,K):
\begin{align}
  \varepsilon_{pp}&=\varepsilon_{0,pp}X^2\rho\,T_6^4, &\varepsilon_{0,pp}&=\val{star-eps0-pp}\,\mathrm{W\,m^3\,kg^{-2}},\\
  \varepsilon_{\rm CNO}&=\varepsilon_{0,\rm CNO}XX_{\rm CNO}\rho\,T_6^{20}, &\varepsilon_{0,\rm CNO}&=\val{star-eps0-cno}\,\mathrm{W\,m^3\,kg^{-2}}.
\end{align}
Reactions happen in the core, $r<0.2R$, which holds a fraction $0.2^3=\val{star-core-mass}$ of the mass.
Its temperature is taken as $T_c$ throughout.
\begin{itemize}
\item $\varepsilon_{pp}=\val{star-eps-pp}$\,W/kg.
\item $\varepsilon_{\rm CNO}=\val{star-eps-cno}$\,W/kg, negligible at this temperature (its $T^{20}$ only takes over above $\sim$18\,MK).
\end{itemize}
The luminosity is then
\begin{equation}
  L_{\rm core}=(\varepsilon_{pp}+\varepsilon_{\rm CNO})\,M_{\rm core}=\val{star-L-core}\,\mathrm{W}=\val{star-L-core-sun}\,L_\odot ,
\end{equation}
which is \val{star-L-ratio} times smaller than the luminosity from question 1.
The surface temperature this luminosity implies,
\begin{equation}
  T=\left(\frac{L_{\rm core}}{4\pi R^2\sigma}\right)^{1/4}=\val{star-T-core-model}\,\mathrm{K},
\end{equation}
is far below the observed \val{star-T}\,K.

\begin{figure}[h]\centering\safegraphic[width=\textwidth]{d2_core}
\caption{Left: temperature profile of the uniform-density star. Right: luminosity of the core against core temperature, for the pp chain and the CNO cycle, compared with the observed luminosity.}\label{fig:core}\end{figure}

\subsection*{Discussion}
The model fails by almost four orders of magnitude.
The culprit is the uniform-density assumption, and \cref{fig:core} (right) shows why it matters so much.
\begin{itemize}
\item \textbf{Density.} A real star is centrally condensed. The Sun's centre is about 100 times denser than its average. A uniform star underestimates the core density, and $\varepsilon\propto\rho$.
\item \textbf{Temperature.} For the same reason it underestimates the core pressure, and so the temperature. With $\varepsilon\propto T^4$ (pp) and $T^{20}$ (CNO), small temperature errors become huge luminosity errors. Matching the observed luminosity at our $\rho_0$ would need $T_c\approx\val{star-T-needed}$\,MK; with the higher real density, less would do.
\item \textbf{Constant core temperature.} Assuming one temperature throughout the core overestimates the outer parts of the core, but this is a smaller effect.
\end{itemize}
Question 1, which uses only the observed radius and temperature, is the reliable value.
The core model shows that the observed luminosity requires a hot, dense, centrally condensed core.

\end{document}
